% Part: normal-modal-logic
% Chapter: axioms-systems
% Section: soundness

\documentclass[../../../include/open-logic-section]{subfiles}

\begin{document}

\olfileid{nml}{prf}{snd}

\olsection{Kasahihan}

Sistem !!^a{derivation} diarani sahih yen kabeh kang bisa
!!{derive}d iku valid. Nalika ngrembug sistem modal, yaiku
!!{derivation}s kang saliyane \Ax{K} uga ngidinake instansi sawetara
!!{formula}s $!A_1$, \dots,~$!A_n$, kita ngarepake saben formula
kang !!{derivable} bener ing saben model kang ndadekake $!A_1$,
\dots, $!A_n$ bener.

\begin{thm}[Teorema Kasahihan]\ollabel{thm:soundness}
  Yen saben instansi saka $!A_1$, \dots, $!A_n$ valid ing
  kelas model $\mClass{C}_1$, \dots, $\mClass{C}_n$
  miturut urutane, mula $\Log{K}!A_1\dots !A_n
  \Proves !B$ ngimplikasikake yen $!B$ valid ing kelas
  model $\mClass{C}_1 \cap \dots \cap \mClass{C}_n$.
\end{thm}

\begin{proof}
  Kanthi induksi marang dawa bukti. Supaya cekak, tetepa $\mClass{C} =
  \mClass{C}_1 \cap \dots \cap \mClass{C}_n$.
  \begin{enumerate}
  \item Basis induksi: Yen $!B$ nduweni bukti kang dawane~$1$, mula
    iku instansi tautologis, instansi
    saka~\Ax{K},\iftag{prvDiamond}{ utawa saka~\Dual{},}{} utawa instansi
    saka salah siji $!A_1$, \dots,~$!A_n$. Ing kasus kapisan, $!B$
    valid ing $\mClass{C}$, amarga instansi tautologis valid ing
    \emph{saben} kelas model, miturut \olref[syn][tau]{prop:valid-taut}.
    Ing kasus kapindho uga mangkono, miturut
    \olref[syn][sch]{prop:Kvalid}\iftag{prvDiamond}{ lan
      \olref[syn][sch]{prop:Dual-valid}}{}. Pungkasan, ing kasus katelu,
    amarga $!B$ valid ing $\mClass{C}_i$ lan $\mClass{C} \subseteq
    \mClass{C}_i$, $!B$ uga valid ing $\mClass{C}$ miturut
    \olref[syn][val]{prop:subset-class}.
  \item Langkah induksi: Upamane $!B$ nduweni bukti kang dawane $k>1$. Yen
    $!B$ iku instansi tautologis, instansi saka~\Ax{K},\iftag{prvDiamond}{ utawa saka~\Dual{},}{} utawa instansi saka salah siji $!A_1$,
    \dots, $!A_n$, kita nindakake kaya ing langkah sadurunge. Dadi upamane $!B$
    dipikolehi kanthi \MP{} saka !!{formula}s sadurunge $!C \lif !B$ lan
    $!C$. Banjur $!C \lif !B$ lan $!C$ nduweni bukti kang dawane $<k$, lan
    miturut hipotesis induksi, loro-lorone valid ing~$\mClass{C}$. Miturut
    \olref[syn][sch]{prop:soundMP}, $!B$ uga valid ing $\mClass{C}$.
    Pungkasan, upamane $!B$ dipikolehi kanthi \Nec{} saka $!C$ (mula
    $!B = \Box!C$). Miturut hipotesis induksi, $!C$ valid ing
    $\mClass{C}$, lan miturut \olref[syn][val]{prop:Nec-rule}, $!B$
    uga valid.
  \end{enumerate}
\end{proof}

\end{document}
